% gz-p7-03-2: 相关系数 r 强度示意——5 个 subplot
\documentclass[tikz,border=4pt]{standalone}
\usepackage{ctex}
\usepackage{amsmath}
\usepackage{pgfplots}
\pgfplotsset{compat=1.18}
\usetikzlibrary{calc}

\begin{document}
\begin{tikzpicture}
% 5 个子图，横向排列
\pgfplotsset{
  small subplot/.style={
    width=4.0cm, height=3.6cm,
    axis lines=center,
    xmin=0, xmax=10, ymin=0, ymax=10,
    xtick=\empty, ytick=\empty,
    tick style={color=black},
    ticklabel style={font=\tiny},
    clip=true,
    title style={font=\small, yshift=-2pt},
  }
}

% r = 1（完美正相关）
\begin{axis}[
  name=r1,
  small subplot,
  title={$r=1$},
]
  \addplot[only marks, mark=*, mark size=1.4pt, blue] coordinates {
    (1,1)(2,2)(3,3)(4,4)(5,5)(6,6)(7,7)(8,8)(9,9)
  };
  \addplot[red, thin, domain=1:9] {x};
\end{axis}

% r = 0.7
\begin{axis}[
  at={($(r1.east)+(1cm,0)$)},
  anchor=west,
  name=r2,
  small subplot,
  title={$r=0.7$},
]
  \addplot[only marks, mark=*, mark size=1.4pt, blue] coordinates {
    (1,2)(2,1.5)(3,4)(4,3)(5,5.5)(6,5)(7,7)(8,6.5)(9,8.5)
    (2,3)(4,4.5)(6,6.5)(8,8)(3,2.5)(7,5.5)
  };
  \addplot[red, thin, domain=1:9] {x*0.85 + 0.7};
\end{axis}

% r = 0
\begin{axis}[
  at={($(r2.east)+(1cm,0)$)},
  anchor=west,
  name=r3,
  small subplot,
  title={$r=0$},
]
  \addplot[only marks, mark=*, mark size=1.4pt, blue] coordinates {
    (2,5)(3,3)(3,7)(5,2)(5,5)(5,8)(7,4)(7,6)(8,2)(8,8)
    (4,4)(6,6)(2,3)(8,5)(4,7)
  };
\end{axis}

% r = -0.7
\begin{axis}[
  at={($(r3.east)+(1cm,0)$)},
  anchor=west,
  name=r4,
  small subplot,
  title={$r=-0.7$},
]
  \addplot[only marks, mark=*, mark size=1.4pt, blue] coordinates {
    (1,8)(2,9)(3,6)(4,7)(5,5.5)(6,4)(7,5)(8,3)(9,1.5)
    (2,7)(4,5.5)(6,5)(8,4.5)(3,8.5)(7,2.5)
  };
  \addplot[red, thin, domain=1:9] {-x*0.85 + 9};
\end{axis}

% r = -1
\begin{axis}[
  at={($(r4.east)+(1cm,0)$)},
  anchor=west,
  name=r5,
  small subplot,
  title={$r=-1$},
]
  \addplot[only marks, mark=*, mark size=1.4pt, blue] coordinates {
    (1,9)(2,8)(3,7)(4,6)(5,5)(6,4)(7,3)(8,2)(9,1)
  };
  \addplot[red, thin, domain=1:9] {-x + 10};
\end{axis}

\node[font=\footnotesize, anchor=north] at ($(r3.south)+(0,-0.6cm)$)
  {从左到右：完美正相关 $\to$ 弱相关 $\to$ 完美负相关};
\end{tikzpicture}
\end{document}
